\section*{Напоминание: задача Коши для уравнения теплопроводности}
\addcontentsline{toc}{section}{Напоминание: уравнение теплопроводности}

Лекция началась с того, что было получено в прошлом семестре. Рассматривалась задача
Коши
\[
  \begin{cases}
    u_t - \alpha u_{xx} = f(x,t), & x \in \R,\ t > 0, \\[2pt]
    u|_{t=0} = \varphi(x),
  \end{cases}
\]
её решение записывается через фундаментальное решение $\Phi$:
\[
  u(x,t) = \intR \Phi(x-y,\,t)\,\varphi(y)\,dy
         + \int\limits_0^t \intR \Phi(x-y,\,t-\tau)\,f(y,\tau)\,dy\,d\tau ,
\]
\[
  \boxed{\ \Phi(x,t) = \frac{1}{2\sqrt{\pi \alpha t}}\;
         e^{-\frac{x^2}{4\alpha t}}\ }
\]

\inputfig{figures/fig-heat-kernel}

\begin{warn}
Рядом в тетради есть плохо читаемая приписка лектора — что-то вроде «как бы то
ни было, это всё приближение; ряд Фурье — это приближение, которое работает на
макроуровне». Смысл, видимо, в том, что само уравнение теплопроводности и его
решение описывают среду усреднённо, а не на уровне отдельных частиц. Стоит
уточнить формулировку.
\end{warn}

\begin{idea}
Структура ответа — свёртка начальных данных с $\Phi$ плюс интеграл Дюамеля по
источнику $f$ — повторится и для волнового уравнения: там роль $\Phi$ сыграет
полусумма сдвигов и интеграл по характеристическому треугольнику.

\medskip
Два свойства $\Phi$ объясняют картинку: $\intR \Phi(x,t)\,dx = 1$ при любом $t>0$
(«тепло никуда не девается»), а высота пика равна
$\Phi(0,t) = \frac{1}{2\sqrt{\pi\alpha t}} \to \infty$ при $t \to 0+$. То есть
$\Phi(\cdot,t) \to \delta$ при $t \to 0+$. При этом $\Phi(x,t) > 0$ сразу для всех
$x$: возмущение в уравнении теплопроводности распространяется мгновенно. Дальше
будет видно, что у волнового уравнения всё наоборот — конечная скорость $a$.
\end{idea}
